\chapter{Background and Related Work}
\label{chap:background}

Building on the clinical motivation established in Chapter 1, this
chapter reviews the technical foundations and related work needed to
motivate QCLand. It focuses on landmark localisation paradigms,
self-supervised vision foundation models, and query-based visual
prediction.


% =================================================================
% 2.1 Fetal Biometry as a Landmark Localisation Task
% =================================================================

\section{Fetal Biometry as a Landmark Localisation Task}
\label{sec:bg-task}

For the purposes of this thesis, fetal biometry is formulated as a
two-landmark localisation problem: given a single 2D ultrasound image
of a standard plane, the task is to predict the two landmark
coordinates that define a specific measurement. This is a direct
localisation task rather than a segmentation task: the required output
is a pair of precise coordinates rather than a delineated anatomical
region or contour.

A characteristic of this formulation is landmark-ordering ambiguity.
The two landmarks defining a diameter measurement have no intrinsic
anatomical first/second identity: the same clinical measurement is
obtained regardless of which landmark is labelled first. Different
datasets may therefore adopt different ordering conventions, and a
model that correctly localises both landmarks but reports them in the
reversed order could be penalised under an order-sensitive evaluation.
This ambiguity must therefore be addressed at training time,
evaluation time, or both \parencite{divece2025multicentre}.


% =================================================================
% 2.2 Deep Learning for Landmark Localisation
% =================================================================

\section{Deep Learning for Landmark Localisation}
\label{sec:bg-landmark-methods}

Common formulations for learning-based 2D landmark localisation
include direct coordinate regression, heatmap regression, and
coordinate classification. Direct coordinate regression collapses
spatial features before predicting numeric $(x,y)$ coordinates;
DeepPose is an early deep-learning example of this formulation
\parencite{toshev2014deeppose}. Although conceptually simple, direct
regression provides less explicit spatial supervision than methods
that retain a structured spatial output, motivating the widespread
use of heatmap- and classification-based formulations.


\subsection{Heatmap Regression}
\label{sec:bg-heatmap}

In heatmap regression, a model produces a dense spatial response map
for each landmark, typically trained to approximate a Gaussian centred
on the ground-truth location
\parencite{tompson2014joint}. A coordinate can then be recovered from
the response map through a discrete maximum with optional local
refinement, or through a differentiable spatial expectation
\parencite{sun2018integral}. The key advantage of this formulation is
that it preserves an explicit correspondence between output locations
and image regions.

Heatmap regression has been developed through progressively refined
convolutional architectures. Stacked hourglass networks
\parencite{newell2016hourglass} introduced repeated
downsampling--upsampling stages with skip connections to combine local
and contextual information. SimpleBaseline
\parencite{xiao2018simple} later showed that a comparatively simple
deconvolutional head on top of a ResNet backbone
\parencite{he2016resnet} can achieve strong keypoint-localisation
performance. HRNet \parencite{sun2019hrnet} took a different approach
by maintaining high-resolution representations throughout the network.
Parallel convolutional streams at multiple resolutions repeatedly
exchange information, allowing the final landmark heatmaps to be
predicted from a high-resolution representation rather than
reconstructed only after a low-resolution bottleneck.

Beyond general pose estimation, heatmap-based localisation has also
been adopted in medical imaging, where explicit spatial modelling can
improve landmark localisation when annotated data are limited
\parencite{payer2019spatial}.


\subsection{HRNet and Landmark Ordering for Fetal Biometry}
\label{sec:bg-hrnet}

The HRNet-W18-based fetal biometry pipeline reported in the
Multicentre Fetal Biometry study
\parencite{divece2025multicentre} adapts high-resolution heatmap
regression to paired fetal landmarks. In addition to the underlying
HRNet architecture, the method incorporates Dynamic Orientation
Determination (DOD) to address landmark-ordering ambiguity. DOD
estimates a measurement-specific direction from the training data and
uses it to impose a consistent ordering on the two landmarks, reducing
errors caused solely by inconsistent annotation conventions.

Related fetal-biometry work has explored direct landmark prediction
using geometric and attention-based constraints
\parencite{gao2023endpoints}.

Together, these studies show how purpose-built landmark architectures
can accommodate the geometric structure of fetal biometry. Their
representations, however, are still learned primarily through
task-specific supervised adaptation, providing a useful contrast to
the large-scale self-supervised representations considered later in
this chapter.


\subsection{Coordinate Classification}
\label{sec:bg-simcc}

Coordinate classification avoids constructing a full two-dimensional
heatmap by discretising the horizontal and vertical axes into
classification bins and predicting a probability distribution over
each coordinate axis. SimCC \parencite{li2022simcc} formulates
keypoint localisation as two independent one-dimensional
classification problems, reducing the quadratic spatial cost of a
full 2D heatmap while permitting fine per-axis discretisation.

RTMPose \parencite{jiang2023rtmpose} builds on this formulation by
combining a CSPNeXt backbone with an RTMCC head based on SimCC
coordinate representations. RTMPose therefore provides a useful
counterpoint to heatmap regression: precise keypoint localisation can
be achieved without maintaining a dense two-dimensional output map.


\subsection{Task-Specific Supervision and Representation Learning}

The distinction between heatmap regression and coordinate
classification lies primarily in their output representations; both
still learn the downstream task through supervised adaptation, even
when initialised from ImageNet-pretrained weights. Their feature
representations are therefore shaped primarily by the labelled
downstream task.

Self-supervised vision foundation models offer a different starting
point. Rather than learning visual representations predominantly from
the downstream annotations, they first acquire general-purpose
representations through large-scale pretraining and are subsequently
adapted to the target task. This shift motivates the central question
for QCLand: whether such representations can be adapted without
sacrificing the spatial precision required for landmark localisation.


% =================================================================
% 2.3 Vision Foundation Models and Query-Based Prediction
% =================================================================

\section{Vision Foundation Models and Query-Based Prediction}
\label{sec:bg-vfm-query}

Two developments are particularly relevant to QCLand:
self-supervised vision transformers, which provide transferable
patch-level visual representations, and query-based prediction
mechanisms, which provide a way to form output-specific
representations through interaction with image features.


\subsection{Vision Transformers}
\label{sec:bg-vit}

A Vision Transformer (ViT, \textcite{dosovitskiy2021vit}) represents
an image as a sequence of fixed-size patches. Each patch is projected
into a token embedding and processed by a stack of transformer encoder
blocks using self-attention. Unlike conventional convolutional
hierarchies, a plain ViT generally maintains a single patch-token
resolution through its encoder, while self-attention allows
information to be exchanged globally across spatial locations.

This property is useful for landmark localisation because each patch
token retains an association with a spatial region while also
incorporating global image context. Features from different
transformer depths therefore share the same spatial grid while
encoding different stages of visual processing, a property that
downstream architectures can exploit.


\subsection{DINOv2 and DINOv3}
\label{sec:bg-dino}

DINOv2 \parencite{oquab2023dinov2} is a self-supervised vision
transformer trained using a teacher--student framework in which the
student is encouraged to produce representations consistent with a
momentum-updated teacher across different augmented views. The
pretraining objective does not require manual semantic annotations.
Earlier work on DINO showed that self-distilled ViTs develop both
strong recognition features and emergent spatial structure
\parencite{caron2021dino}. DINOv2 scales this approach to large,
curated image collections, producing general-purpose visual
representations that transfer effectively to downstream tasks such as
classification, segmentation, and depth estimation.

Of particular relevance to landmark localisation are DINOv2's dense
patch-level representations. Individual patch tokens remain spatially
associated with image regions while encoding increasingly contextual
visual information through the transformer layers. Such
representations provide a richer pretrained starting point for
spatial prediction than learning all visual features solely from a
small task-specific labelled dataset.

DINOv3 \parencite{simeoni2025dinov3} further scales the DINO family
and introduces changes intended to preserve the quality of dense
representations during large-scale training. Gram anchoring
regularises patch-level feature relationships during training, while
rotary position embeddings replace the learned absolute position
embeddings used by DINOv2. DINOv2 and DINOv3 therefore represent two
successive generations of self-supervised vision foundation models
that provide strong dense visual representations for downstream
adaptation.


\subsection{Query-Based Prediction: From DETR to EoMT}
\label{sec:bg-eomt}

Learned query representations have become an established mechanism for
structured visual prediction. DETR \parencite{carion2020detr}
introduced learned object queries that interact with image features
through a transformer decoder, with bipartite Hungarian matching used
to associate predictions with ground-truth objects. Mask2Former
\parencite{cheng2022mask2former} extended the query-based formulation
to segmentation, where each query predicts both a semantic category
and a spatial mask.

The Encoder-only Mask Transformer (EoMT,
\textcite{kerssies2025eomt}) simplifies this architecture by removing
the separate transformer decoder. Instead, learned mask-query tokens
are inserted directly into the ViT patch-token sequence partway
through the encoder. Subsequent transformer blocks allow queries and
image tokens to interact through the backbone's own self-attention,
so the ViT itself performs the query--image information exchange.

For spatial prediction, each EoMT query is transformed and compared
with transformed dense patch features through a per-location dot
product,

\begin{equation}
    \label{eq:bg-einsum}
    M_q(x,y)
    =
    f(q)^\top g\bigl(z(x,y)\bigr),
\end{equation}

where $q$ denotes a query representation, $z(x,y)$ is the visual
feature at spatial position $(x,y)$, and $f$ and $g$ are learned
transformations. This formulation is well matched to region-mask
prediction, but segmentation performance alone does not establish
its suitability for sharply localised landmark targets.

EoMT also applies supervision at multiple transformer depths. The
predicted masks can influence subsequent query-to-image attention,
while Hungarian matching associates a variable number of
segmentation instances with the available queries. These mechanisms
reflect the setting for which EoMT was designed:
variable-cardinality region segmentation rather than fixed-count
landmark localisation.


% =================================================================
% 2.4 Summary and Research Gap
% =================================================================

\section{Summary and Research Gap}
\label{sec:bg-summary}

The literature reviewed above brings together three complementary
lines of work. Purpose-built landmark architectures such as HRNet
provide high-resolution spatial localisation and have already been
adapted to fetal biometry, but rely primarily on task-specific
supervised training. Self-supervised vision foundation models such as
DINOv2 and DINOv3 provide general-purpose, spatially grounded patch
representations learned through large-scale pretraining, but do not
themselves produce landmark coordinates. Query-based architectures
such as EoMT allow learned queries to form output-specific
representations through interaction with ViT image tokens, but were
developed for variable-cardinality segmentation rather than precise
paired-landmark prediction.

Fetal biometry combines these requirements in a different way: each
measurement is defined by a fixed pair of landmarks, their ordering
may be ambiguous, and accurate measurement depends on spatially
precise coordinate recovery. The resulting adaptation problem is
therefore to combine query-conditioned foundation-model
representations with a spatial readout suited to precise landmark
localisation. \Cref{chap:method} introduces QCLand to address this
gap.
